\chapter{Related Work}
\label{chapter:related-work}

This chapter reviews existing approaches to improving the effectiveness of
fuzzing for regression testing, with a particular focus on seed selection and
generation. We first examine traditional regression fuzzing techniques that
adapt the fuzzing process to prioritize changed or affected code, followed by
approaches for selecting, reusing, and generating seed corpora. We then review
recent LLM-driven seed-generation techniques and distinguish them by their scope
of targeting and their integration with the fuzzing process. Finally, we
synthesize these approaches into a comparative taxonomy and identify the
research gap our work addresses.

\section{Traditional Regression Fuzzing}
\label{sec:rw-regression-fuzzing}

Regression fuzzing aims to accelerate the discovery of bugs introduced by recent
code changes by prioritizing the exploration of commit-related program regions.
Existing approaches have predominantly pursued this goal by adapting the fuzzing
engine through target-aware instrumentation, guidance mechanisms, and scheduling
strategies.

The seminal work in this area is AFLGO~\cite{bohmeDirectedGreyboxFuzzing2017},
which introduced \ac{DGF}. AFLGO instruments the target binary with distance
metrics derived from the control flow graph and uses a simulated annealing power
schedule to prioritize seeds that execute closer to specified target locations.
This established the foundation for commit-aware fuzzing by steering execution
toward modified code.

Subsequent work has focused on reducing the overhead of directed fuzzing to
better suit Continuous Integration (CI/CD) environments.
CIDFuzz~\cite{zhangCIDFuzzFuzzTesting2023} uses differential analysis to
identify change points and dynamically allocate fuzzing resources, while
Pazzer~\cite{sharmaEffectiveFuzzingCI2024a} replaces expensive intra-procedural
CFG analysis with incremental call-graph construction. Similarly,
SelectFuzz~\cite{luoSelectFuzzEfficientDirected2023} reduces runtime overhead by
instrumenting only data- and control-dependent paths relevant to the target.

While early \ac{DGF} approaches targeted only modified lines, later work
recognized that regression bugs frequently manifest in code indirectly affected
by a commit. WAFLGO~\cite{xiangCriticalCodeGuided2024} therefore extends the
target scope beyond the immediate code diff by incorporating transitively
affected program regions calculated via data-flow dependency analysis into its
guidance strategy. 

Overall, traditional regression fuzzing has primarily focused on adapting the
fuzzing engine through target-aware analysis, instrumentation, and execution
guidance. Since our work addresses the same objective of accelerating regression
bug discovery in commit-affected code, these approaches provide the relevant
regression-fuzzing context for our work. Unlike these methods, however, we
investigate ways to improve the initial seed corpus rather than modify the
fuzzing process itself.

\section{Corpus Selection, Reuse, and Augmentation}
\label{sec:rw-corpus}

The quality of the initial seed corpus substantially influences fuzzing
effectiveness, which is particularly relevant under the limited execution
budgets of CI/CD environments. Empirical studies have shown that seed selection
can significantly affect fuzzing
performance~\cite{herreraSeedSelectionSuccessful2021}, while corpus distillation
can reduce redundant inputs and improve the efficiency of subsequent
fuzzing~\cite{herreraCorpusDistillationEffective2020}.

Consequently, prior work has investigated the reuse and selection of inputs from existing fuzzing campaigns.
Continuous fuzzing approaches, aimed at regression-oriented fuzzing, retain and select corpora across commit-triggered campaigns, based on functional coverage~\cite{yooImprovingConfigurabilityUnitlevel2021,kimBugOssBenchmarkRealworld2024}, or distance to the change location~\cite{canakciTargetFuzzUsingDARTs2022}.

These approaches motivate treating the corpus as a persistent artifact rather than reconstructing it for every campaign.
Our approach, therefore, retains the existing corpus, distills it to the most relevant seeds, and augments it with newly generated seeds that target not only the parts changed by the current commit but also the transitively affected parts.
This combines previously acquired program coverage with commit-specific inputs while allowing the resulting corpus to be used by an otherwise unmodified fuzzer.

\section{Traditional Seed Generation}
\label{sec:rw-traditional-generation}

While regression fuzzing has historically focused on modifying execution
guidance and scheduling within the fuzzing engine, parallel research has
investigated how to create a high-quality initial seed corpus to give the
fuzzing engine a substantial head start.

Early input seed generation relied heavily on manual, domain-specific generation
scripts, such as Domato~\cite{GoogleprojectzeroDomato2026} and
Freedom~\cite{xuFREEDOMEngineeringStateoftheArt2020a} for HTML files. To
minimize manual engineering, subsequent work leveraged formal grammar and format
specifications. Tools like
FormatFuzzer~\cite{dutraFormatFuzzerEffectiveFuzzing2023} compile binary format
specifications into executable generators and mutators, whereas frameworks such
as Fandango~\cite{zamudioamayaFANDANGOEvolvingLanguageBased2025,
bettscheiderCombiningInputConstraints2026} synthesize inputs during fuzzing by
enforcing grammar constraints aligned with execution goals.

To eliminate the need for manual grammar specification, deep learning approaches
introduced data-driven seed synthesis by training machine learning models
directly on existing valid seed corpora. Early probabilistic models like
Skyfire~\cite{wangSkyfireDataDrivenSeed2017} learned input distributions to
generate seeds for complex parsers, while
LearnAndFuzz~\cite{godefroidLearnFuzzMachineLearning2017} utilized Recurrent
Neural Networks for sequence-based input synthesis. Subsequent work, such as
Smartseed~\cite{lyuSmartSeedSmartSeed2019}, employed Generative Adversarial
Networks to generate structurally valid seed files.

Despite their effectiveness, traditional generation techniques remain
constrained: grammar-based methods require substantial domain expertise to
formalize specifications, while classical ML approaches struggle to generalize
to unseen program semantics without extensive retraining data. These limitations
paved the way for \ac{LLM} approaches capable of zero-shot semantic reasoning
across diverse codebases.

\section{LLM-Driven Seed Generation}
\label{sec:rw-llm-generation}

Recent research has increasingly adopted Large Language Models and autonomous
agents to overcome the generality bottlenecks of traditional seed generation. We
categorize them by their \emph{targeting scope and context awareness}: broad, non-targeted approaches; single-point-targeted approaches; and commit-derived, multi-point-targeted regression approaches.

\subsection{Broad- and Non-Targeted Approaches}
\label{subsec:rw-broad}

Broad \ac{LLM} approaches leverage generative models to maximize the overall program
coverage or to synthesize syntactically valid inputs without focusing on a
specific vulnerability or code diff.

Early efforts, such as OSS-Fuzz
Gen~\cite{OssfuzzgenLlm_toolkitCorpus_generatorpy}, used one-shot prompting,
providing the model with project descriptions and highly relevant source code
snippets to synthesize initial corpus files. Building upon this,
SeedMind~\cite{shiHarnessingLargeLanguage2024} employs a more structured
methodology, traversing the Dynamic Call Graph in a breadth-first manner and
utilizing coverage-guided feedback to iteratively prompt the \ac{LLM} to write
and refine seed-generation scripts to maximize overall code coverage.

Other techniques integrate \ac{LLM} generation directly into the dynamic fuzzing
loop rather than focusing solely on the initial corpus.
G2Fuzz~\cite{EnablingLowCostComprehensive} generates initial seeds by tasking
the \ac{LLM} to write Python scripts that, in turn, create specific input
formats (like PDFs). Another \ac{LLM} is then dynamically activated to generate
mutation-like candidate inputs when the underlying fuzzer becomes locally stuck.
Similarly, Fuzz4All~\cite{xiaFuzz4AllUniversalFuzzing2024} processes
human-provided documentation and code snippets to generate initial seeds using
an auto-prompting technique, then scores them based on input validity. The
prompt is then forwarded to the main fuzzing loop, where the \ac{LLM}
continuously alternates between input generation, mutation, and semantic
equivalence transformation. Fuzz4All operates without a dynamic feedback
mechanism, terminating either when a crash is detected or when its computational
budget expires.

\subsection{Single-Point-Targeted Approaches}
\label{subsec:rw-point-targeted}

Recognizing that broad generation often fails to reach deep or complex program
logic, a substantial line of work steers \acp{LLM} toward a single known
location, typically a reported vulnerability. These approaches differ mainly in
how much context they assemble and how the generated inputs are consumed.

ISC4DGF~\cite{xuISC4DGFEnhancingDirected2024} supplies the model with project
metadata, driver source, CVE descriptions, and patch diffs, and synthesizes
Python generation scripts whose outputs seed a subsequent fuzzing campaign.
HGFuzzer~\cite{xuDirectedGreyboxFuzzing2025} extends this to an all-in-one
solution that additionally generates harnesses and infers execution conditions
via a secondary model instance.

SeedSmith~\cite{zhuSeedSmithLLMDrivenSeed2026} replaces static context assembly
with an agent that actively queries the codebase for function bodies,
caller--callee relationships, and type definitions, refining its generation
scripts against execution feedback before handing the corpus to a fuzzer.
RandLuzz~\cite{fengFuzzingRandomnessReasoning2025} goes further still, building a
knowledge base from CVE summaries and bug reports, synthesizing inputs that step
along the call chain toward the target, and finally producing a specialized
mutator that operates inside the fuzzing loop.

Although several of these approaches consume patch diffs as context, their
generation is steered toward a single sink location rather than the change as a
whole; we therefore classify them as point-targeted. Such methods presuppose
that the target location is known in advance, which does not hold when a commit
modifies multiple locations, and the faulty one has not yet been identified.

\subsection{Commit-Derived, Multi-Point-Targeted Approaches}
\label{subsec:rw-commit-aware}

While point-targeted approaches excel at reaching single static locations,
regression testing requires evaluating distributed changes spread across an
entire commit. The primary approach addressing this setting is
Cleverest~\cite{liuEvaluatingLLMBasedRegression2026}.

Cleverest passes the raw commit diff together with a project description to an
\ac{LLM} and tasks it with generating an input that triggers a bug introduced by
the commit. Each candidate is executed, and the outcome, whether the changed
code was reached, whether the program output differed, and whether a bug was
triggered, is recorded in an attempt history that is fed back into the next
prompt. This loop continues until a bug is found or the budget is exhausted. If
no bug is triggered directly, the inputs accumulated during the process are
handed to a fuzzer as a starting corpus, giving it a head start toward the bug.

Two aspects of this design bound its scope. Its targets are drawn from the raw
diff, so code that a change affects only transitively, through data- or
control-dependent paths, is not represented in the generation context. Its
corpus is also constructed anew for each commit, so coverage accumulated by
prior campaigns does not inform the result.

\section{Comparative Taxonomy and Research Gap}
\label{sec:rw-gap}
\Cref{tab:fuzzing_taxonomy} summarizes the related work along the dimensions introduced above: the targeting scope of generated seeds; whether generation completes before the fuzzing campaign or is invoked from within it; whether directly changed and transitively affected code is targeted; and whether an existing corpus is retained and extended.
Two gaps emerge:
\begin{description}
  \item[Transitive scope is unused outside the engine.]
    On the engine side, WAFLGO extends its guidance to code that a change affects through data-flow dependencies.
    On the corpus side, no generation-based approach draws on more than the diff itself, and approaches that distill an existing corpus do not consider transitively affected code either.
  \item[Augmentation is absent throughout.]
    Approaches that generate seeds construct them anew for each campaign, while the work that does treat corpora as persistent artifacts~\cite{herreraSeedSelectionSuccessful2021, yooImprovingConfigurabilityUnitlevel2021, canakciTargetFuzzUsingDARTs2022, liImprovingSeedQuality2025} selects from what already exists rather than generating new inputs for a specific change.
\end{description}

\begin{tcolorbox}[colback=black!4, colframe=black!40, boxrule=0.4pt,
                  left=6pt, right=6pt, top=6pt, bottom=6pt]
To the best of our knowledge, neither has a seed corpus been \textbf{distilled toward the code a commit affects} beyond the lines it changes, nor has an existing corpus been \textbf{augmented by an LLM agent} with inputs targeting that code, and the two steps have never been \textbf{combined in a single pipeline}.
\end{tcolorbox}

Our approach closes these gaps, relative to Cleverest~\cite{liuEvaluatingLLMBasedRegression2026}, our closest prior work, and the only other generation-based approach targeting commits.
Where Cleverest derives its targets from the raw diff, we derive target regions from the \ac{CoI} of the changed statements, capturing code a commit affects indirectly through data or control dependencies.
And where Cleverest constructs a corpus from scratch, we first distill the existing corpus into seeds that exercise these regions, then augment it with inputs generated by a \ac{LLM} for the remaining uncovered parts.
Moreover, where Cleverest relies on iterative prompting, in which the model receives feedback only through successive prompts, we employ an \ac{LLM} agent that autonomously manages its context by exploring the codebase, the existing corpus, and coverage feedback to refine its inputs.

% The following three commands belong in the preamble.
\newcommand{\yes}{\ding{51}}
\newcommand{\no}{\ding{55}}
\newcommand{\na}{\textendash}

\begin{table}[t]
\centering
\footnotesize
\setlength{\tabcolsep}{5pt}
\renewcommand{\arraystretch}{1.25}
\begin{tabularx}{\linewidth}{@{}>{\raggedright\arraybackslash}X
  >{\centering\arraybackslash}p{0.13\linewidth}
  *{4}{>{\centering\arraybackslash}p{0.085\linewidth}}@{}}
\toprule
\textbf{Approach} & \textbf{Scope} & \textbf{Pre-Fuzz} &
\textbf{Commit} & \textbf{Affected} & \textbf{Augment} \\
\midrule

\multicolumn{6}{@{}l}{\itshape Regression fuzzing without seed generation}\\[2pt]
AFLGo~\cite{bohmeDirectedGreyboxFuzzing2017},
CIDFuzz~\cite{zhangCIDFuzzFuzzTesting2023},
Pazzer~\cite{sharmaEffectiveFuzzingCI2024a},
SelectFuzz~\cite{luoSelectFuzzEfficientDirected2023}
  & None & \na & \yes & \no  & \no \\
WAFLGO~\cite{xiangCriticalCodeGuided2024}
  & None & \na & \yes & \yes & \no \\

\midrule
\multicolumn{6}{@{}l}{\itshape Seed generation without LLMs}\\[2pt]
FormatFuzzer~\cite{dutraFormatFuzzerEffectiveFuzzing2023},
Fandango~\cite{zamudioamayaFANDANGOEvolvingLanguageBased2025,
bettscheiderCombiningInputConstraints2026}
  & Broad & \no & \no & \no & \no \\
Domato~\cite{GoogleprojectzeroDomato2026},
FREEDOM~\cite{xuFREEDOMEngineeringStateoftheArt2020a},
Skyfire~\cite{wangSkyfireDataDrivenSeed2017},
LearnAndFuzz~\cite{godefroidLearnFuzzMachineLearning2017},
SmartSeed~\cite{lyuSmartSeedSmartSeed2019}
  & Broad & \yes & \no & \no & \no \\

\midrule
\multicolumn{6}{@{}l}{\itshape LLM-based seed generation}\\[2pt]
G2Fuzz~\cite{EnablingLowCostComprehensive},
Fuzz4All~\cite{xiaFuzz4AllUniversalFuzzing2024}
  & Broad  & \no  & \no  & \no  & \no \\
OSS-Fuzz-Gen~\cite{OssfuzzgenLlm_toolkitCorpus_generatorpy},
SeedMind~\cite{shiHarnessingLargeLanguage2024}
  & Broad  & \yes & \no  & \no  & \no \\
ISC4DGF~\cite{xuISC4DGFEnhancingDirected2024},
HGFuzzer~\cite{xuDirectedGreyboxFuzzing2025},
RandLuzz~\cite{fengFuzzingRandomnessReasoning2025}
  & Single & \no  & \no  & \no  & \no \\
SeedSmith~\cite{zhuSeedSmithLLMDrivenSeed2026}
  & Single & \yes & \no  & \no  & \no \\
Cleverest~\cite{liuEvaluatingLLMBasedRegression2026}
  & Multi  & \yes & \yes & \no  & \no \\

\midrule
\textbf{Ours} & \textbf{Multi} & \yes & \yes & \yes & \yes \\
\bottomrule
\end{tabularx}

\caption{Taxonomy of fuzzing and seed-generation approaches by covered domain.}
\label{tab:fuzzing_taxonomy}
\end{table}